\documentclass[10pt,a4paper]{article}
\usepackage[a4paper,left=21mm,right=21mm,top=17mm,bottom=19mm]{geometry}
\usepackage{fontspec}
\usepackage{enumitem}
\usepackage{hyperref}
\usepackage{tabularx}
\usepackage{ragged2e}
\usepackage{fancyhdr}
\setmainfont{Helvetica}
\hypersetup{pdftitle={Johannes Roth - Machine Learning Engineer},pdfauthor={Johannes Roth},pdfsubject={Machine learning engineering, applied data science, and research},colorlinks=true,urlcolor=black}
\pagestyle{fancy}
\fancyhf{}
\renewcommand{\headrulewidth}{0pt}
\cfoot{\fontsize{8.5}{10}\selectfont Johannes Roth \quad | \quad \thepage\ / 2}
\setlength{\parindent}{0pt}
\setlength{\parskip}{0pt}
\setlength{\RaggedRightRightskip}{0pt plus 1fil}
\setlist[itemize]{leftmargin=12pt,itemsep=3pt,topsep=4pt,parsep=0pt,partopsep=0pt}
\newcommand{\sectiontitle}[1]{\par\vspace{12pt}{\fontsize{12}{14}\selectfont\bfseries\MakeUppercase{#1}}\par\vspace{3pt}\hrule height 0.4pt\vspace{4pt}}
\newcommand{\role}[3]{\par\vspace{8pt}\begin{tabularx}{\textwidth}{@{}Xr@{}}\textbf{#1}&{\small #2}\\\end{tabularx}\par\textit{#3}\par}
\begin{document}
\begin{center}
{\fontsize{26}{30}\selectfont\bfseries Johannes Roth}\par
\vspace{5pt}
{\fontsize{12}{14}\selectfont Machine Learning Engineer}\par
\vspace{7pt}
{\small \href{mailto:johannes@roth24.de}{johannes@roth24.de}\quad
\href{https://linkedin.com/in/johannes-roth-a67baa141}{LinkedIn}\quad
\href{https://github.com/andropar}{GitHub}\quad
\href{https://scholar.google.com/citations?user=U3Z3J2gAAAAJ}{Scholar}\quad
\href{https://jroth.space}{jroth.space}\quad Leipzig, Germany}\par
\vspace{4pt}
{\small Available from January 2027}
\end{center}
\fontsize{10.5}{13}\selectfont
\RaggedRight
\hyphenpenalty=10000
\exhyphenpenalty=10000
Machine learning engineer with experience taking ML and data projects from initial requirements to deployment and operation. Background in production computer vision, large-scale data processing, and independent research in computational neuroscience.

\sectiontitle{Experience}
\begin{samepage}
\role{PhD Researcher - Computational Cognitive Neuroscience}{May 2022 -- Present}{Max Planck Institute for Human Cognitive and Brain Sciences \& University of Giessen}
\begin{itemize}
\item Conducted independent research at the intersection of \textbf{machine learning and computational neuroscience}, focusing on vision models and their relationship to human brain activity.
\item Developed methods for model evaluation and experimental design, combining computational analyses with fMRI studies.
\item Developed \textbf{GPU-based ML pipelines} for dataset curation and feature extraction, and released \href{https://huggingface.co/datasets/andropar/relaion2b-natural}{large-scale image datasets} for computer-vision research.
\item Independently developed and maintained research software and cloud-based data services for \href{https://laion-fmri.hebartlab.com}{\textbf{LAION-fMRI}} and \href{https://re-vision-initiative.org/}{\textbf{re:vision}}, including public websites and annotation tools.
\item Supervised working students and interns and collaborated with researchers across computer science and neuroscience.
\end{itemize}
\end{samepage}

\begin{samepage}
\role{Research Assistant - Machine Learning in Medicine}{Jun 2021 -- May 2022}{ScaDS.AI Dresden/Leipzig}
\begin{itemize}
\item Developed and evaluated deep-learning models for \textbf{medical image analysis}, including brain-tumor segmentation.
\item Built predictive models for clinical data, with a focus on mortality prediction and uncertainty estimation.
\end{itemize}
\end{samepage}

\begin{samepage}
\role{Data Scientist / ML Engineer (Working Student)}{Nov 2019 -- May 2021}{CHECK24, Travel Vertical}
\begin{itemize}
\item Independently designed and deployed a \textbf{production computer-vision service} for \textbf{20 million hotel images}, covering model development and backend engineering.
\item Prototyped additional computer-vision applications for image retrieval, duplicate detection, and image-quality assessment.
\item Optimized recommendation-data pipelines through incremental processing, reducing update time by over \textbf{95\%}.
\end{itemize}
\end{samepage}

\begin{samepage}
\role{Full-Stack Developer (Freelance)}{Oct 2020 -- May 2021}{Kimetric UG}
\begin{itemize}
\item Designed and delivered Django-based websites for research organizations, including \href{https://hebartlab.com}{\textbf{Hebart Lab}} and the \href{https://things-initiative.org}{\textbf{THINGS Initiative}}.
\item Gathered client requirements and managed deployment, Linux hosting, and ongoing application maintenance.
\end{itemize}
\end{samepage}

\newpage
\begin{samepage}
\role{Data Scientist (Working Student)}{Oct 2018 -- Sep 2019}{Webdata Solutions GmbH (now Vistex)}
\begin{itemize}
\item Developed machine-learning models for \textbf{product matching and image classification} in e-commerce.
\item Improved image-based matching accuracy from \textbf{12\% to 83\%} against a classification baseline in an internal cross-retailer evaluation.
\end{itemize}
\end{samepage}

\begin{samepage}
\role{Data Analyst (Working Student)}{Feb 2018 -- Aug 2018}{Mercateo Services GmbH \& Co. KG}
\begin{itemize}
\item Developed data-integration, reporting, and testing workflows for an \textbf{AWS-based data warehouse} using Apache NiFi and Pentaho.
\end{itemize}
\end{samepage}

\sectiontitle{Technical Skills}
\begin{itemize}
\item \textbf{ML \& Analysis:} Python, SQL, PyTorch, TensorFlow, scikit-learn; computer vision, embeddings, model evaluation, statistical analysis, experimental design.
\item \textbf{Software \& Data:} FastAPI, Flask, Django, Redis, Docker, Git, CI/CD; data pipelines, similarity search, dataset curation.
\item \textbf{Cloud \& Compute:} AWS (S3, IAM), Linux, SLURM/HPC, multi-GPU processing.
\item \textbf{Languages:} German (native), English (fluent).
\end{itemize}

\sectiontitle{Open Source}
\begin{itemize}
\item Contributed to \href{https://github.com/ViCCo-Group/thingsvision}{\textbf{thingsvision}}, an open-source library for computer-vision research, with work on software architecture, model integration, and memory efficiency.
\end{itemize}

\sectiontitle{Education}
\role{Doctoral Candidate (Dr. rer. nat.)}{May 2022 -- Present}{University of Giessen}
Research focus: machine learning and computational neuroscience.

\role{M.Sc. Computer Science}{2017 -- 2021}{Leipzig University; grade 1.2}
\begin{itemize}
\item \textbf{Master's thesis (grade 1.1):} Applied generative deep learning to study visual processing in the human brain.
\end{itemize}

\role{B.Sc. Business Information Systems}{2014 -- 2017}{Leipzig University; grade 1.5}

\sectiontitle{Selected Publications}
\begin{itemize}
\item F. P. Mahner*, \textbf{J. Roth*}, et al. \href{https://arxiv.org/abs/2605.13675}{\textit{Characterizing Universal Object Representations Across Vision Models}}. \textbf{NeurIPS 2026} (accepted). Shared first author.
\item \textbf{J. Roth*}, Y. Duan*, et al. \href{https://doi.org/10.1038/s44271-025-00236-3}{\textit{Ten principles for reliable, efficient, and adaptable coding in psychology and cognitive neuroscience}}. Communications Psychology, 2025. Shared first author.
\item \textbf{J. Roth} et al. \href{https://doi.org/10.1007/978-3-031-08999-2_23}{\textit{Multi-plane UNet++ Ensemble for Glioblastoma Segmentation}}. MICCAI BrainLes 2021 workshop proceedings, published 2022.
\end{itemize}

\sectiontitle{Recognition \& Service}
\begin{itemize}
\item Presented research at \href{https://openreview.net/forum?id=T9k6KkZoca}{\textbf{CCN 2025 (oral)}}.
\item Elected representative for \textbf{180+ doctoral researchers} at MPI CBS (2023--2024).
\item \textbf{CMBB Replication Award (2025)} for contributions to reliable coding practices in neuroscience.
\end{itemize}
\end{document}
